% ============================================================
%  TD Chapitre 2 : Logique et algebre de Boole
% ============================================================
\clearpage
\chapter{Logique et algèbre de Boole}
\chaptermark{Logique et algèbre de Boole}

\begin{exercice}[du français aux propositions]
Dans un laboratoire informatique de l'institut, on considère les deux propositions :
\[
P : \text{« la connexion Internet fonctionne »} \quad ; \quad Q : \text{« le groupe électrogène est en marche »}.
\]
Traduire en phrases françaises simples les six propositions suivantes :
\[
\lnot P \qquad ; \qquad P \land Q \qquad ; \qquad P \lor Q \qquad ; \qquad Q \lor \lnot P \qquad ; \qquad \lnot P \lor \lnot Q \qquad ; \qquad \lnot (P \land Q).
\]
Comparer la signification des deux dernières et indiquer, sans table, si elles doivent être équivalentes.
\end{exercice}

\begin{exercice}[quatre tables de vérité]
Construire, sur les quatre combinaisons de valeurs de $(P, Q)$ prises dans l'ordre du comptage binaire, les tables de vérité des quatre expressions suivantes, en donnant une colonne par sous-formule :
\[
\lnot P \land Q \qquad ; \qquad P \lor \lnot Q \qquad ; \qquad \lnot (P \land Q) \qquad ; \qquad \lnot P \lor \lnot Q.
\]
Comparer les colonnes de $\lnot (P \land Q)$ et de $\lnot P \lor \lnot Q$, puis conclure sur l'équivalence annoncée par la première loi de De Morgan.
\end{exercice}

\begin{exercice}[équivalences, réciproque et contraposée]
\begin{enumerate}
\item Dresser une table unique de $P$, $Q$, $\lnot (P \lor Q)$ et $\lnot P \land \lnot Q$ ; comparer les deux dernières colonnes et énoncer la loi vérifiée.
\item Dresser une table de $P$, $Q$, $\lnot P$, $P \rightarrow Q$ et $\lnot P \lor Q$ ; conclure.
\item Dresser une table de $P$, $Q$, $P \rightarrow Q$, $Q \rightarrow P$ et $\lnot Q \rightarrow \lnot P$. Identifier la réciproque et la contraposée de $P \rightarrow Q$ parmi les deux dernières colonnes et dire lesquelles sont équivalentes à $P \rightarrow Q$.
\item En déduire la forme correcte de la négation de l'énoncé « si la sauvegarde est terminée, alors l'ordinateur s'éteint » parmi les deux phrases suivantes, en la justifiant par la table : (a) « la sauvegarde est terminée et l'ordinateur ne s'éteint pas » ; (b) « si la sauvegarde n'est pas terminée, alors l'ordinateur ne s'éteint pas ».
\end{enumerate}
\end{exercice}

\begin{exercice}[simplification d'expressions booléennes]
Simplifier les trois expressions suivantes en citant chaque loi utilisée, $A$, $B$ et $C$ étant des variables booléennes :
\[
X = A \cdot B + A \cdot (B + C) \qquad ; \qquad Y = (A + B) \cdot (A + C) \qquad ; \qquad Z = A \cdot (A + B).
\]
Vérifier les trois simplifications par une table de vérité unique sur les huit combinaisons de $(A, B, C)$, qui juxtapose les colonnes de $X$ et de sa forme simplifiée, de $Y$ et de sa forme simplifiée, de $Z$ et de $A$.
\end{exercice}

\begin{exercice}[opérations bit à bit et chiffrement par le ou exclusif]
\begin{enumerate}
\item On donne les octets $P = 01100111$ et $Q = 11110001$. Calculer bit à bit $P \land Q$, $P \lor Q$ et $P \veebar Q$.
\item Un message de quatre bits $M = 1011$ est chiffré par le masque $K = 0110$ : calculer le message chiffré $C = M \veebar K$, puis retrouver $M$ en calculant $C \veebar K$. Justifier ce résultat par les propriétés du ou exclusif.
\item Un mini-circuit utilise trois portes : la sortie de la porte ET des entrées $A$ et $B$ et la sortie de la porte ET de $\overline{A}$ et $B$ alimentent ensemble une porte OU. Écrire l'expression booléenne de la sortie, la simplifier et dresser la table de vérité qui confirme la simplification.
\end{enumerate}
\end{exercice}
